% GENERATED FILE. Edit canonical lessons, then run npm run generate:books.
% canonical-source-hash: fnv1a64:14b1ba5cbc0c9c98
% canonical-lessons: BN-C303-sedin, BN-C303-otit, BN-C303-perono, BN-C303-kalokrome, BN-C303-smoron-kora

\chapter{Long Ago}
\label{ch:bn-long-ago}
\hlchaptermodality{303}

\begin{chapteropening}
\textbf{By the end of this chapter:} \emph{I can place events in the past: that day, long ago, as the years went by.}

Complete the last lesson of chapter 303: I can place events in the past: that day, long ago, as the years went by.
\end{chapteropening}

\section[\texorpdfstring{\bn{সেদিন} (sedin)}{সেদিন (sedin)}]{\bn{সেদিন} (sedin) — that day; the other day}

\label{lesson:BN-C303-sedin}



\begin{quote}
Before the new one: say the Bengali for to bring about, then the Bengali for gradually.
\end{quote}

\subsection*{You'll want to know: \bn{সেদিন}}
\textbf{\bn{সেদিন}} — \emph{sedin} — \textquotedblleft{}that day; the other day\textquotedblright{}.

\begin{quote}
\bn{সেদিন} \bn{বিকেলে} \bn{হঠাৎ} \bn{ঝড়} \bn{এল।}
\end{quote}

\begin{quote}
\emph{sedin bikele hôṭhāt jhôṛ elô.}
\end{quote}

\begin{quote}
That day a storm blew up suddenly in the late afternoon.
\end{quote}

\begin{cousinweb}
\textbf{\bn{সেদিন}} (\emph{sedin}) joins \bn{সে}, 'that', to \textbf{\bn{দিন}}, a day — Sanskrit \emph{dina} (CDIAL 6328) — into one word: that day, the day the story is about. In talk it also means 'the other day', a few days ago.
\end{cousinweb}

\begin{grammarlens}[title={What you've built}]
One more word for telling what happened.
\end{grammarlens}

\subsection*{Your turn}
\begin{itemize}
\raggedright
  \item \emph{Say it:} \emph{sedin}
  \item \emph{Say it:} \emph{sedin}, once more
  \item \emph{Say it:} \emph{sedin bikele hôṭhāt jhôṛ elô.} — That day a storm blew up suddenly in the late afternoon.
  \item \emph{Recall:} say the Bengali for to bring about, then the Bengali for gradually, then say \emph{sedin} again
\end{itemize}

\begin{tcolorbox}[breakable,colback=teal!4,colframe=teal!35!black,title={Before you move on}]
What does \bn{সেদিন} mean? (\textquotedblleft{}That day; the other day\textquotedblright{}.) Say it once more.
\end{tcolorbox}
\section[\texorpdfstring{\bn{অতীত} (ôtīt)}{অতীত (ôtīt)}]{\bn{অতীত} (ôtīt) — the past}

\label{lesson:BN-C303-otit}



\begin{quote}
Before the new one: say the Bengali for gradually, then the Bengali for that day.
\end{quote}

\subsection*{You'll want to know: \bn{অতীত}}
\textbf{\bn{অতীত}} — \emph{ôtīt} — \textquotedblleft{}the past\textquotedblright{}.

\begin{quote}
\bn{দাদু} \bn{প্রায়ই} \bn{অতীতের} \bn{গল্প} \bn{বলেন।}
\end{quote}

\begin{quote}
\emph{dādu prāyi ôtīter gôlpo bôlen.}
\end{quote}

\begin{quote}
Grandfather often talks about the past.
\end{quote}

\begin{cousinweb}
\textbf{\bn{অতীত}} (\emph{ôtīt}) is Sanskrit \emph{atīta}, 'gone by, past' (Monier-Williams), from \emph{ati-}, 'beyond', and \emph{√i}, 'to go'.
\end{cousinweb}

\begin{grammarlens}[title={What you've built}]
One more word for telling what happened.
\end{grammarlens}

\subsection*{Your turn}
\begin{itemize}
\raggedright
  \item \emph{Say it:} \emph{ôtīt}
  \item \emph{Say it:} \emph{ôtīt}, once more
  \item \emph{Say it:} \emph{dādu prāyi ôtīter gôlpo bôlen.} — Grandfather often talks about the past.
  \item \emph{Recall:} say the Bengali for gradually, then the Bengali for that day, then say \emph{ôtīt} again
\end{itemize}

\begin{tcolorbox}[breakable,colback=teal!4,colframe=teal!35!black,title={Before you move on}]
What does \bn{অতীত} mean? (\textquotedblleft{}The past\textquotedblright{}.) Say it once more.
\end{tcolorbox}
\section[\texorpdfstring{\bn{পেরোনো} (perono)}{পেরোনো (perono)}]{\bn{পেরোনো} (perono) — to pass by; to go past}

\label{lesson:BN-C303-perono}



\begin{quote}
Before the new one: say the Bengali for that day, then the Bengali for the past.
\end{quote}

\subsection*{You'll want to know: \bn{পেরোনো}}
\textbf{\bn{পেরোনো}} — \emph{perono} — \textquotedblleft{}to pass by; to go past\textquotedblright{}.

\begin{quote}
\bn{বছর} \bn{পেরোল}, \bn{কিন্তু} \bn{রাজকুমার} \bn{ফিরল} \bn{না।}
\end{quote}

\begin{quote}
\emph{bôchhôr perolô, kintu rājkumār phirlô nā.}
\end{quote}

\begin{quote}
Years went by, but the prince did not return.
\end{quote}

Past (he/she …): \bn{সেই} \bn{নৌকা} \bn{নদী} \bn{পেরোল}, \emph{sei noukā nôdī perolô}.

\begin{cousinweb}
\textbf{\bn{পেরোনো}} (\emph{perono}) is usually explained as built on \bn{পার}, 'the far bank' — Sanskrit \emph{pāra} (CDIAL 8100) — with the verb ending: to get to the far side of a river, or of a stretch of time.
\end{cousinweb}

\begin{grammarlens}[title={What you've built}]
One more word for telling what happened.
\end{grammarlens}

\subsection*{Your turn}
\begin{itemize}
\raggedright
  \item \emph{Say it:} \emph{perono}
  \item \emph{Say it:} \emph{perono}, once more
  \item \emph{Say it:} \emph{bôchhôr perolô, kintu rājkumār phirlô nā.} — Years went by, but the prince did not return.
  \item \emph{Recall:} say the Bengali for that day, then the Bengali for the past, then say \emph{perono} again
\end{itemize}

\begin{tcolorbox}[breakable,colback=teal!4,colframe=teal!35!black,title={Before you move on}]
What does \bn{পেরোনো} mean? (\textquotedblleft{}To pass by; to go past\textquotedblright{}.) Say it once more.
\end{tcolorbox}
\section[\texorpdfstring{\bn{কালক্রমে} (kālôkrôme)}{কালক্রমে (kālôkrôme)}]{\bn{কালক্রমে} (kālôkrôme) — in the course of time}

\label{lesson:BN-C303-kalokrome}



\begin{quote}
Before the new one: say the Bengali for the past, then the Bengali for to go past.
\end{quote}

\subsection*{You'll want to know: \bn{কালক্রমে}}
\textbf{\bn{কালক্রমে}} — \emph{kālôkrôme} — \textquotedblleft{}in the course of time\textquotedblright{}.

\begin{quote}
\bn{কালক্রমে} \bn{গ্রামটা} \bn{একটা} \bn{বড়} \bn{শহর} \bn{হল।}
\end{quote}

\begin{quote}
\emph{kālôkrôme grāmṭā ekṭā bôṛo shôhôr hôlô.}
\end{quote}

\begin{quote}
In the course of time the village became a big town.
\end{quote}

\begin{cousinweb}
\textbf{\bn{কালক্রমে}} (\emph{kālôkrôme}) is Sanskrit \emph{kālakrama}, 'the lapse of time' (Monier-Williams), in the -e form: \emph{kāla}, 'time', and \emph{krama}, 'a step'.
\end{cousinweb}

\begin{grammarlens}[title={What you've built}]
One more word for telling what happened.
\end{grammarlens}

\subsection*{Your turn}
\begin{itemize}
\raggedright
  \item \emph{Say it:} \emph{kālôkrôme}
  \item \emph{Say it:} \emph{kālôkrôme}, once more
  \item \emph{Say it:} \emph{kālôkrôme grāmṭā ekṭā bôṛo shôhôr hôlô.} — In the course of time the village became a big town.
  \item \emph{Recall:} say the Bengali for the past, then the Bengali for to go past, then say \emph{kālôkrôme} again
\end{itemize}

\begin{tcolorbox}[breakable,colback=teal!4,colframe=teal!35!black,title={Before you move on}]
What does \bn{কালক্রমে} mean? (\textquotedblleft{}In the course of time\textquotedblright{}.) Say it once more.
\end{tcolorbox}
\section[\texorpdfstring{\bn{স্মরণ} \bn{করা} (smôrôn kôrā)}{স্মরণ করা (smôrôn kôrā)}]{\bn{স্মরণ} \bn{করা} (smôrôn kôrā) — to commemorate; to call to mind}

\label{lesson:BN-C303-smoron-kora}



\begin{quote}
Before the new one: say the Bengali for to go past, then the Bengali for in the course of time.
\end{quote}

\subsection*{You'll want to know: \bn{স্মরণ} \bn{করা}}
\textbf{\bn{স্মরণ} \bn{করা}} — \emph{smôrôn kôrā} — \textquotedblleft{}to commemorate; to call to mind\textquotedblright{}.

\begin{quote}
\bn{আজ} \bn{শহর} \bn{সেই} \bn{ভূমিকম্পের} \bn{দিনটা} \bn{স্মরণ} \bn{করল।}
\end{quote}

\begin{quote}
\emph{āj shôhôr sei bhūmikômper dinṭā smôrôn kôrlô.}
\end{quote}

\begin{quote}
Today the city commemorated the day of that earthquake.
\end{quote}

Past (he/she …): \bn{তিনি} \bn{বাবাকে} \bn{স্মরণ} \bn{করলেন}, \emph{tini bābāke smôrôn kôrlen}.

\begin{cousinweb}
\textbf{\bn{স্মরণ} \bn{করা}} (\emph{smôrôn kôrā}) puts Sanskrit \emph{smaraṇa}, 'the act of remembering' (Monier-Williams), before \bn{করা}. It is more formal than \textbf{\bn{মনে} \bn{রাখা}} and is the word for honouring a memory in public; \textbf{\bn{স্মৃতি}} comes from the same root.
\end{cousinweb}

\begin{grammarlens}[title={What you've built}]
One more word for telling what happened.
\end{grammarlens}

\subsection*{Your turn}
\begin{itemize}
\raggedright
  \item \emph{Say it:} \emph{smôrôn kôrā}
  \item \emph{Say it:} \emph{smôrôn kôrā}, once more
  \item \emph{Say it:} \emph{āj shôhôr sei bhūmikômper dinṭā smôrôn kôrlô.} — Today the city commemorated the day of that earthquake.
  \item \emph{Recall:} say the Bengali for to go past, then the Bengali for in the course of time, then say \emph{smôrôn kôrā} again
\end{itemize}

\begin{tcolorbox}[breakable,colback=teal!4,colframe=teal!35!black,title={Before you move on}]
What does \bn{স্মরণ} \bn{করা} mean? (\textquotedblleft{}To commemorate; to call to mind\textquotedblright{}.) Say it once more.
\end{tcolorbox}
